% GENERATED FILE. Edit canonical lessons, then run npm run generate:books.
% canonical-source-hash: fnv1a64:5bfb54412218b111
% canonical-lessons: TA-C103-grandson, TA-C103-granddaughter, TA-C103-neighbour, TA-C103-familyname, TA-C103-petname

\chapter{Grandchildren, Neighbours, Names}
\label{ch:ta-grandchildren-names}
\hlchaptermodality{103}

\begin{chapteropening}
\textbf{By the end of this chapter:} \emph{I can name a grandson and granddaughter, a neighbour, a family name and a pet name.}

Complete the last lesson of chapter 103: i can name a grandson and granddaughter, a neighbour, a family name and a pet name.
\end{chapteropening}

\section[pēraṉ]{\ta{பேரன்} (pēraṉ) — a grandson}

\label{lesson:TA-C103-grandson}



\begin{quote}
Before the new one: say the Tamil for a daughter-in-law, then the Tamil for a son-in-law.
\end{quote}

\subsection*{You'll want to know: \ta{பேரன்}}
\textbf{\ta{பேரன்}} (\emph{pēraṉ}) — \textquotedblleft{}a grandson\textquotedblright{}.

\textbf{-\ta{அன்}} marks a man, as in \textbf{\ta{மருமகன்}}.

\begin{grammarlens}[title={What you've built}]
A grandson.
\end{grammarlens}

\subsection*{Your turn}
\begin{itemize}
\raggedright
  \item \emph{Say it:} \emph{pēraṉ}
  \item \emph{Say it:} \emph{pēraṉ}, once more
  \item \emph{Say it:} say \emph{eṉ pēraṉ}
  \item \emph{Recall:} say the Tamil for a daughter-in-law, then the Tamil for a son-in-law, then say \emph{pēraṉ} again
\end{itemize}

\begin{tcolorbox}[breakable,colback=teal!4,colframe=teal!35!black,title={Before you move on}]
What does \ta{பேரன்} mean? (\textquotedblleft{}A grandson\textquotedblright{}.) Say it once more.
\end{tcolorbox}
\section[pētti]{\ta{பேத்தி} (pētti) — a granddaughter}

\label{lesson:TA-C103-granddaughter}



\begin{quote}
Before the new one: say the Tamil for a son-in-law, then the Tamil for a grandson.
\end{quote}

\subsection*{You'll want to know: \ta{பேத்தி}}
\textbf{\ta{பேத்தி}} (\emph{pētti}) — \textquotedblleft{}a granddaughter\textquotedblright{}.

It pairs with \textbf{\ta{பேரன்}} the way \textbf{\ta{பாட்டி}} pairs with \textbf{\ta{தாத்தா}}.

\begin{grammarlens}[title={What you've built}]
Grandson and granddaughter.
\end{grammarlens}

\subsection*{Your turn}
\begin{itemize}
\raggedright
  \item \emph{Say it:} \emph{pētti}
  \item \emph{Say it:} \emph{pētti}, once more
  \item \emph{Say it:} say \emph{eṉ pēraṉ, eṉ pētti}
  \item \emph{Recall:} say the Tamil for a son-in-law, then the Tamil for a grandson, then say \emph{pētti} again
\end{itemize}

\begin{tcolorbox}[breakable,colback=teal!4,colframe=teal!35!black,title={Before you move on}]
What does \ta{பேத்தி} mean? (\textquotedblleft{}A granddaughter\textquotedblright{}.) Say it once more.
\end{tcolorbox}
\section[pakkattu vīṭṭukkārar]{\ta{பக்கத்து} \ta{வீட்டுக்காரர்} (pakkattu vīṭṭukkārar) — a neighbour}

\label{lesson:TA-C103-neighbour}



\begin{quote}
Before the new one: say the Tamil for a grandson, then the Tamil for a granddaughter.
\end{quote}

\subsection*{You'll want to know: \ta{பக்கத்து} \ta{வீட்டுக்காரர்}}
\textbf{\ta{பக்கத்து} \ta{வீட்டுக்காரர்}} (\emph{pakkattu vīṭṭukkārar}) — \textquotedblleft{}a neighbour\textquotedblright{}, literally \textquotedblleft{}the person of the house beside\textquotedblright{}.

\textbf{\ta{பக்கம்}} is \textquotedblleft{}side\textquotedblright{}; \textbf{\ta{வீடு}} is \textquotedblleft{}house\textquotedblright{}.

\begin{grammarlens}[title={What you've built}]
A neighbour, spelled out as the house beside.
\end{grammarlens}

\subsection*{Your turn}
\begin{itemize}
\raggedright
  \item \emph{Say it:} \emph{pakkattu vīṭṭukkārar}
  \item \emph{Say it:} \emph{pakkattu vīṭṭukkārar}, once more
  \item \emph{Say it:} say \emph{pakkattu vīṭṭukkārar}
  \item \emph{Recall:} say the Tamil for a grandson, then the Tamil for a granddaughter, then say \emph{pakkattu vīṭṭukkārar} again
\end{itemize}

\begin{tcolorbox}[breakable,colback=teal!4,colframe=teal!35!black,title={Before you move on}]
What does \ta{பக்கத்து} \ta{வீட்டுக்காரர்} mean? (\textquotedblleft{}A neighbour\textquotedblright{}.) Say it once more.
\end{tcolorbox}
\section[kuṭumpap peyar]{\ta{குடும்பப்} \ta{பெயர்} (kuṭumpap peyar) — a family name}

\label{lesson:TA-C103-familyname}



\begin{quote}
Before the new one: say the Tamil for a granddaughter, then the Tamil for a neighbour.
\end{quote}

\subsection*{You'll want to know: \ta{குடும்பப்} \ta{பெயர்}}
\textbf{\ta{குடும்பப்} \ta{பெயர்}} (\emph{kuṭumpap peyar}) — \textquotedblleft{}a family name\textquotedblright{}.

\textbf{\ta{பெயர்}} is the \textquotedblleft{}name\textquotedblright{} of your first introductions; \textbf{\ta{குடும்பம்}} is \textquotedblleft{}family\textquotedblright{}. Many Tamil names use a father's initial where English uses a surname.

\begin{grammarlens}[title={What you've built}]
A family name, and what often stands in for it.
\end{grammarlens}

\subsection*{Your turn}
\begin{itemize}
\raggedright
  \item \emph{Say it:} \emph{kuṭumpap peyar}
  \item \emph{Say it:} \emph{kuṭumpap peyar}, once more
  \item \emph{Say it:} say \emph{kuṭumpap peyar}
  \item \emph{Recall:} say the Tamil for a granddaughter, then the Tamil for a neighbour, then say \emph{kuṭumpap peyar} again
\end{itemize}

\begin{tcolorbox}[breakable,colback=teal!4,colframe=teal!35!black,title={Before you move on}]
What does \ta{குடும்பப்} \ta{பெயர்} mean? (\textquotedblleft{}A family name\textquotedblright{}.) Say it once more.
\end{tcolorbox}
\section[cellappeyar]{\ta{செல்லப்பெயர்} (cellappeyar) — a pet name, a nickname}

\label{lesson:TA-C103-petname}



\begin{quote}
Before the new one: say the Tamil for a neighbour, then the Tamil for a family name.
\end{quote}

\subsection*{You'll want to know: \ta{செல்லப்பெயர்}}
\textbf{\ta{செல்லப்பெயர்}} (\emph{cellappeyar}) — \textquotedblleft{}a pet name\textquotedblright{}, literally \textquotedblleft{}darling-name\textquotedblright{}.

\textbf{\ta{செல்லம்}} is \textquotedblleft{}darling, pet\textquotedblright{}.

\begin{grammarlens}[title={What you've built}]
Grandchildren, a neighbour, a family name, a pet name.
\end{grammarlens}

\subsection*{Your turn}
\begin{itemize}
\raggedright
  \item \emph{Say it:} \emph{cellappeyar}
  \item \emph{Say it:} \emph{cellappeyar}, once more
  \item \emph{Say it:} say \emph{eṉ cellappeyar ...}
  \item \emph{Recall:} say the Tamil for a neighbour, then the Tamil for a family name, then say \emph{cellappeyar} again
\end{itemize}

\begin{tcolorbox}[breakable,colback=teal!4,colframe=teal!35!black,title={Before you move on}]
What does \ta{செல்லப்பெயர்} mean? (\textquotedblleft{}A pet name, a nickname\textquotedblright{}.) Say it once more.
\end{tcolorbox}
